
\documentclass[11pt]{article}

\usepackage[utf8]{inputenc}
\usepackage[T1]{fontenc}
\usepackage[spanish,english]{babel}
\usepackage{amsmath,amssymb,amsthm}
\usepackage[hidelinks]{hyperref}
\usepackage[margin=2.5cm]{geometry}
\usepackage{enumitem}

\setlength{\parskip}{0.6em}
\setlength{\parindent}{0pt}

\newtheorem{proposition}{Proposition}
\newtheorem{hypothesis}{Hypothesis}

\title{Behavioral Model of Vote Buying and Vote Selling}
\date{\today}

\begin{document}

\maketitle

\section{Objectivos del nuevo modelo}

\begin{itemize}

\item \textbf{Formalizar explicitamente el mecanismo de overspending por aversion a las perdidas.}  
El modelo nuevo no puede quedarse en compensacion minima tipo $\Delta$ si el paper quiere sostener que, cuando los partidos inician, tienden a sobrepagar para asegurarse contra la derrota. Ese mecanismo debe salir del modelo, no solo de la narrativa.

\item \textbf{Introducir una referencia clara de perdidas y ganancias para los partidos.}  
Si la intuicion central es prospect theory (que se basa en un punto de referencia), el modelo necesita un punto de referencia frente al cual la derrota electoral sea codificada como perdida, y la victoria o la no-accion no queden tratadas simetricamente. Sin eso, no hay fundamento formal para la asimetria entre vote buying y vote selling.

\item \textbf{Hacer que la secuencia de movimientos afecte realmente la activacion del ``loss frame''.}  
La diferencia entre party-initiated y voter-initiated no deberia ser solo quien mueve primero en un juego estandar, sino cuando el partido entra en una situacion psicologica o estrategica en la que percibe el riesgo electoral como perdida. Eso tiene que estar modelado, no solo descrito.

\item \textbf{Derivar formalmente H3, no dejarla como una implicacion solo verbal.}  
El nuevo modelo deberia mostrar con claridad bajo que condiciones los votantes obtienen mayores payoffs cuando los partidos inician el intercambio, y por que los partidos quedan relativamente mejor cuando los votantes inician. Esa comparacion de surplus tiene que salir como resultado formal.

\item \textbf{Mantener y conservar H1 y H2 como resultados formales del modelo.}  
El modelo no debe perder lo que ya funciona en el paper (se refiere al modelo que esta en el paper mismo):
\begin{enumerate}
\item core targeting cuando inician los partidos,
\item venta a partidos electoralmente fuertes cuando inician los votantes.
\end{enumerate}
El nuevo modelo debe seguir generando estos resultados, pero ahora ademas incorporar bien la logica de surplus y overspending (como en nuestro paper pasado, overspending ocurre cuando van ganando).

\item \textbf{Asegurar que H2 salga como resultado condicional, no universal.}  
El modelo debe dejar explicito que vender al partido electoralmente fuerte y distante ocurre solo cuando su stake electoral es suficientemente alto para sostener un acuerdo factible. Esa condicion tiene que quedar transparente en el equilibrio.

\item \textbf{Modelar aceptacion/rechazo de precios de forma consistente con el experimento.}  
En el juego de vote selling, si el votante pide un precio y el partido acepta, eso debe implicar una transaccion implementada (como ocurre en el experimento). El modelo tiene que reflejar que los precios pedidos son reservation prices vinculantes y no cheap talk.

\item \textbf{Conectar mejor el modelo con el hecho de que en el experimento no siempre hay pivotality exacta.}  

Idealmente, el nuevo modelo deberia o bien:
\begin{enumerate}
\item incorporar explicitamente estados con votante pivotal y no pivotal, o
\item dejar clarisimo que el caso pivotal es el benchmark formal y que lo demas es una extension.
\end{enumerate}

Ahora mismo en el modelo que esta en el paper, esa relacion sigue siendo mas debil de lo ideal.

\item \textbf{Permitir que electoral strength entre como stake real del partido, no solo como etiqueta.}  
La fuerza electoral debe afectar la disposicion a pagar de manera estructural, no solo intuitiva. Tiene que verse en la funcion objetivo del partido o en las condiciones de equilibrio.

\item \textbf{Hacer visible por que un partido fuerte pero ideologicamente distante puede ser el mejor comprador para el votante.}  
Esa logica esta en el paper, pero el nuevo modelo debe mostrarla con mayor limpieza: el votante compara distancia ideologica versus capacidad/disposicion del partido a pagar, y el segundo puede dominar bajo ciertas condiciones.

\item \textbf{Formalizar mejor el trade-off entre compensacion ideologica minima y sobrepago conductual.}  
Seria ideal que el modelo muestre dos niveles distintos: uno ``baseline'' de compensacion minima por distancia ideologica, y otro adicional generado por aversion a la perdida o sobreaseguramiento electoral. Asi el overspending no reemplaza toda la logica anterior, sino que la extiende.

\item \textbf{Evitar que el modelo implique solo minimal compensation en el caso central.}  
Ese fue el principal problema de consistencia. Si el equilibrio principal vuelve a dar solo $\Delta$, entonces el paper otra vez queda tensionado con la interpretacion empirica y con la agenda previa de ustedes sobre overspending.

\item \textbf{Dejar muy claro que parte del modelo es estandar y que parte es behavioral.}  
Si se combina un juego estrategico clasico con un componente prospect-theoretic, eso debe quedar ordenado y transparente. El lector tiene que poder ver exactamente que supuesto nuevo genera la diferencia entre vote buying y vote selling.

\item \textbf{Generar predicciones comparativas directamente observables en el experimento.}  
El modelo nuevo deberia producir predicciones sobre: ofertas, precios pedidos, aceptacion/rechazo, y payoffs. Mientras mas directamente mapee a las variables observadas, mejor queda la conexion teoria-empiria.

\item \textbf{Lograr una arquitectura donde abstract, teoria, modelo y resultados digan exactamente lo mismo.}  
El modelo reescrito deberia permitir afirmar sin matices raros: que el modelo formal explica targeting, explica bajo que condiciones los votantes venden al partido fuerte, y explica por que la iniciativa redistribuye surplus. Ahora esa ultima parte todavia no esta suficientemente formalizada.

\end{itemize}

What comes next it's an extension of the baseline model in the paper, which incorporates a behavioral mechanism that distinguishes party-initiated vote buying from voter-initiated vote selling. The key intuition is that the sequencing of moves does not merely redistribute bargaining power; it also determines whether parties evaluate the transaction in a loss frame. When parties initiate the exchange, they confront the possibility of electoral defeat at the moment of commitment and may therefore over-insure against that outcome. When voters initiate the exchange, by contrast, parties can screen proposals before entering the same commitment frame, which limits voters' ability to extract surplus. This behavioral asymmetry generates distinct predictions about targeting, pricing, and payoffs across institutional variants.


\section{Propuesta de nuevo modelo que se supone logra los objectivos de arriba}


\paragraph{Players, preferences, and electoral stakes.}

Consider two parties, $i\in\{A,B\}$, and one voter $j$. The policy space is one-dimensional. Party $i$ is located at $\gamma_i$, and the voter has ideal point $x_j$. If party $i$ wins, the voter receives ideological utility

\begin{equation}
u_j(\gamma_i)=D-\lvert x_j-\gamma_i\rvert,
\end{equation}

where $D$ is large enough to ensure non-negative utility levels. Let

\begin{equation}
i^\ast=\arg\max_{i\in\{A,B\}} u_j(\gamma_i)
\end{equation}

denote the voter's ideologically preferred party. Define the voter's ideological advantage from supporting the core party as

\begin{equation}
\Delta\equiv u_j(\gamma_{i^\ast})-u_j(\gamma_{-i^\ast})>0.
\end{equation}

This is the minimal ideological compensation required to induce the voter to support the less-preferred party.

If the voter supports party $i$ and receives transfer $t_i\geq 0$, then total voter utility is

\begin{equation}
U_j(i,t_i)=u_j(\gamma_i)+t_i.
\end{equation}

Each party values winning the election at $W_i>0$. Let $\pi_i\in(0,1)$ denote party $i$'s subjective probability that the voter's support is decisive for avoiding defeat. The instrumental electoral stake of securing the vote is therefore

\begin{equation}
R_i\equiv \pi_i W_i.
\end{equation}

Higher values of $R_i$ capture greater electoral risk or greater strategic value from obtaining the voter's support.

\paragraph{Prospect-theoretic party preferences.}

Unlike the baseline model, parties do not evaluate outcomes linearly. Instead, party $i$ evaluates outcomes relative to a reference point $r_i$, interpreted as the payoff the party expects if it does not secure the voter's support under the prevailing electoral environment. Let the party's value function be

\begin{equation}
V_i(z-r_i)=
\begin{cases}
z-r_i & \text{if } z\geq r_i,\\
-\lambda_i(r_i-z) & \text{if } z<r_i,
\end{cases}
\end{equation}

where $\lambda_i>1$ captures loss aversion. Outcomes below the reference point are thus weighted more heavily than equally sized gains above it.

\subsection*{Game 1: Party-Initiated Vote Buying}

In the vote-buying game, party $i$ chooses a transfer offer $s_i\geq 0$ before the voter chooses which party to support. If party $i$ secures the vote, its monetary payoff is $W_i-s_i$; if it fails to secure the vote, its payoff is $0$.

The distinctive behavioral assumption is that, once a party actively decides whether to insure against defeat, the relevant reference point becomes avoiding the electorally salient loss of failing to secure pivotal support. Let $L_i>0$ denote the perceived electoral loss associated with that failure once the party has entered the decision frame. Let $p_i(s_i)$ denote the probability that offer $s_i$ secures the vote, where $p_i'(s_i)>0$. Party $i$ solves

\begin{equation}
\max_{s_i\geq 0}\; p_i(s_i)(W_i-s_i)+(1-p_i(s_i))[-\lambda_i L_i].
\end{equation}
The first-order condition is
\begin{equation}
p_i'(s_i)\bigl(W_i-s_i+\lambda_i L_i\bigr)=p_i(s_i).
\end{equation}

Relative to a standard expected-payoff benchmark, loss aversion increases the marginal benefit of raising the offer because $\lambda_iL_i$ enters positively in the left-hand side. As a result, whenever $\lambda_i>1$, the optimal offer $s_i^{VB}$ may exceed the minimal compensating transfer $\Delta$. In other words, the core voter remains cheaper to buy in ideological terms, but party initiative may induce overspending because parties insure against a prospectively salient electoral loss.

This yields the model's first implication. In party-initiated vote buying, ideological proximity continues to structure targeting, so transfers are directed primarily toward core voters. Yet, conditional on targeting those voters, parties may offer more than the minimal compensating transfer required to overturn ideology.

\subsection*{Game 2: Voter-Initiated Vote Selling}

In the vote-selling game, the voter first proposes reservation prices $(a_A,a_B)$, and each party then decides whether to accept. If party $i$ accepts and is chosen, its payoff is $W_i-a_i$; if it rejects, it simply forgoes the transaction.

The critical behavioral difference is that the party evaluates the proposal before entering a commitment frame. Rejection is therefore not coded as a realized loss relative to an insured reference point. Let the reference point in this stage be the status quo of non-transaction, normalized to $r_i=0$. Then party $i$ accepts if and only if

\begin{equation}
V_i(W_i-a_i)\geq V_i(0).
\end{equation}
Because acceptance yields a gain only when $a_i\leq W_i$, while rejection leaves the party at the status quo, the acceptance condition reduces to
\begin{equation}
a_i\leq W_i.
\end{equation}

Thus, in the vote-selling game, loss aversion does not inflate willingness to pay in the same way. Parties can reject costly proposals without entering the same loss frame that operates when they themselves initiate the exchange.

The voter chooses reservation prices to maximize utility subject to the parties' acceptance constraints. Suppose party $k$ has the highest electoral stake. The voter solves

\begin{equation}
\max_{a_A,a_B}\; \max\{u_j(\gamma_A)+a_A,\;u_j(\gamma_B)+a_B\}
\quad\text{s.t.}\quad
a_i\leq W_i \;\; \forall i.
\end{equation}

If $k\neq i^\ast$, the voter sells to the electorally stronger but ideologically distant party whenever

\begin{equation}
u_j(\gamma_k)+a_k\geq u_j(\gamma_{i^\ast})+a_{i^\ast},
\end{equation}
which is equivalent to
\begin{equation}
a_k-a_{i^\ast}\geq \Delta.
\end{equation}

A sufficient condition for such an equilibrium to exist is that the stronger party's feasible payment advantage is large enough to offset ideological distance, namely

\begin{equation}
W_k-W_{i^\ast}\geq \Delta.
\end{equation}

This yields the second implication of the model. When voters initiate the exchange, selling is reallocated toward electorally stronger parties whenever their feasible payment advantage is sufficiently large to compensate for ideological disadvantage.

\paragraph{Payoff implication.}

The third implication concerns surplus extraction. When parties initiate vote buying, loss aversion operates at the moment of commitment and may induce over-insurance against defeat. When voters initiate vote selling, parties screen requests before entering the same behavioral frame and can reject costly offers without incurring the same perceived loss. Under $\lambda_i>1$ and sufficiently salient electoral losses $L_i$, expected transfers satisfy

\begin{equation}
\mathbb{E}[s_i^{VB}]>\mathbb{E}[a_i^{VS,\text{accepted baseline}}],
\end{equation}

implying that voters extract greater surplus under party-initiated vote buying than under voter-initiated vote selling.

Taken together, the behavioral model delivers three predictions. First, party initiative preserves core targeting but may generate overspending relative to the minimal compensating transfer. Second, voter initiative reallocates selling toward electorally stronger parties when their feasible payment advantage offsets ideological distance. Third, voters obtain higher expected payoffs when parties, rather than voters, initiate the exchange.

\section{Antiguo modelo: el modelo que esta en el paper}

\paragraph{Players, preferences, and electoral stakes.}

There are two parties, $i \in \{A,B\}$, and a single pivotal voter $j$. The policy space is one-dimensional, $\gamma \in \Gamma = \{1,\dots,100\}$, with $\gamma_A < \gamma_B$. The voter has an ideal point $x_j$ drawn independently from the same distribution. If party $i$ wins, the voter receives ideological utility
\begin{equation}
u_j(\gamma_i) \;=\; D - \lvert x_j - \gamma_i \rvert,
\end{equation}
with $D$ large enough to keep utilities non-negative. Let
\begin{equation}
i^\ast \;=\; \arg\max_{i \in \{A,B\}} u_j(\gamma_i)
\end{equation}
denote the voter's preferred (core) party.

The voter is pivotal with probability $\pi>0$. Party $i$ values winning at $W_i>0$, so the expected benefit of securing the pivotal vote is:
\begin{equation}
R_i \;=\; \pi W_i.
\end{equation}
We interpret $R_i$ as party $i$'s electoral stake or risk.

Transfers consist of non-negative private benefits. Transfers offered by parties in the vote-buying game are denoted $s_i$, and minimum acceptable transfers requested by the voter in the vote-selling game are denoted $a_i$.

The voter's total utility from voting for party $i$ and receiving transfer $t_i$ is:
\begin{equation}
U_j(i,t_i) \;=\; u_j(\gamma_i) + t_i.
\end{equation}

%----------------------------------------------------------
\subsection*{Game 1: Party-Initiated Vote Buying}

\begin{enumerate}
\item Nature draws $(x_j,\gamma_A,\gamma_B)$ and $(R_A,R_B)$ and reveals them to all players.
\item Each party $i$ simultaneously chooses a transfer offer $s_i \in [0,\bar{B}]$.
\item The voter observes $(s_A,s_B)$ and chooses which party to support.
\end{enumerate}

If the voter selects party $i$, she receives $s_i$, and party $i$ receives net payoff $R_i - s_i$. The losing party receives zero.

Define the ideological advantage of the core party as:
\begin{equation}
\Delta \;=\; u_j(\gamma_{i^\ast}) - u_j(\gamma_{-i^\ast}) \;>\; 0.
\end{equation}

This is the minimal compensating transfer needed to overturn the voter's ideological preference.

In symmetric environments with $R_A = R_B = R$ and $R > \Delta$, standard undercutting arguments imply:
\begin{equation}
s^{VB}_{i^\ast} = \Delta, \qquad s^{VB}_{-i^\ast} = 0.
\end{equation}

Thus the core party buys the vote at minimal cost.

\begin{proposition}[Core Targeting in Vote Buying]\label{prop:core_buying}
In the party-initiated vote-buying game, whenever purchasing the pivotal vote is profitable, there exists an equilibrium in which the ideologically preferred party $i^\ast$ buys the vote with the minimal compensating transfer $\Delta$, while the opponent does not buy. Transfers therefore concentrate on core voters.
\end{proposition}

%----------------------------------------------------------
\subsection*{Game 2: Voter-Initiated Vote Selling}

\begin{enumerate}
\item Nature draws $(x_j,\gamma_A,\gamma_B)$ and $(R_A,R_B)$ and reveals them.
\item The voter proposes a pair of minimum acceptable transfers $(a_A,a_B)$.
\item Each party $i$ accepts ($b_i = 1$) or rejects ($b_i = 0$).
\end{enumerate}

Party $i$ accepts if and only if
\begin{equation}
a_i \;<\; R_i,
\end{equation}
so $R_i$ is party $i$'s maximum willingness to pay. Equivalently, the voter can set requests arbitrarily close to each party's willingness to pay by choosing $a_i = R_i - \varepsilon$ for some $\varepsilon>0$.

Let $k \in \{A,B\}$ denote the electorally stronger party, so $R_k > R_{-k}$, and consider the empirically relevant case where $k \neq i^\ast$. (Here, $W_i$ denotes party $i$'s utility of winning, whereas $k$ indexes the party with higher electoral stakes $R_i$.)

If both parties accept, the voter chooses $k$ whenever
\begin{equation}
u_j(\gamma_k) + a_k \;\ge\; u_j(\gamma_{i^\ast}) + a_{i^\ast}
\quad\Longleftrightarrow\quad
a_k - a_{i^\ast} \;\ge\; \Delta.
\end{equation}

There exists a band of prices satisfying
\begin{equation}
\Delta \;\le\; a_k - a_{i^\ast} \;\le\; R_k - R_{i^\ast},
\end{equation}
within which both parties accept but the vote is sold to $k$.

The voter maximizes her material payoff by setting
\begin{equation}
a_k^{VS} = R_k - \varepsilon, \qquad \varepsilon > 0,
\end{equation}
with $a_{i^\ast}$ chosen low enough to satisfy the inequality above.

\begin{proposition}[Selling to the Opponent Winning Party]\label{prop:selling_strong}
If the electorally stronger party $k$ is ideologically distant ($k \neq i^\ast$ and $R_k > R_{i^\ast}$), then in the vote-selling game there exists an equilibrium in which both parties accept the voter's proposals but the vote is sold to $k$. Vote selling thus directs transfers toward the opponent expected to win the election.\footnote{This result is conditional as it arises only when the stronger party's electoral stake is sufficiently large to support a feasible agreement despite ideological distance.}
\end{proposition}

\end{document}